% Generated from locale/tr/content/set-theory/ord-arithmetic/exponentiation.tex; do not hand-edit.


\olfileid[tr]{sth}{ord-arithmetic}{expo}
\olsection{Sıra Sayılarında Üs Alma}

Şimdi sıra sayılarında üs almaya geçiyoruz. Ne yazık ki sıra sayılarında üs
alma için \emph{zarif} bir sentetik tanım yoktur.

Elbette açık sentetik tanımlar \emph{vardır}. İşte bunlardan biri.
$\text{finfun}(\alpha,\beta)$, $f\colon\beta\to\alpha$ olup
$\Setabs{\gamma\in\beta}{f(\gamma)\neq 0}$ kümesi bir doğal sayıyla eş
güçlü olan bütün $f$ fonksiyonlarının kümesi olsun. $f\sqsubset g$ ancak ve
ancak $f\neq g$ ve
$\gamma_0=\text{max}\Setabs{\gamma\in\beta}{f(\gamma)\neq g(\gamma)}$
için $f(\gamma_0)<g(\gamma_0)$ olacak biçimde
$\text{finfun}(\alpha,\beta)$ üzerinde bir iyi sıralama tanımlayın. O zaman
$\ordexpo{\alpha}{\beta}$'yı
$\ordtype{\text{finfun}(\alpha,\beta),\sqsubset}$ olarak tanımlayabiliriz.
Potter bu açık tanımı kullanır ve hemen ardından şöyle açıklar:
\begin{quote}
	Bu sıralamanın seçimi, yalnızca uygun özyineleme koşuluna uyan bir sıra
	sayısı üs alma tanımı elde etme isteğimizce belirlenir\ldots{} ve onu
	gözde canlandırmak, sıralı toplamı ya da sıralı çarpımı
	gözde canlandırmaktan çok daha zordur.
	\citep[p.~199]{Potter2004}
\end{quote}
Aynen öyle. Toplamayı ``$\alpha$'nın bir kopyasının ardından $\beta$'nın
bir kopyası''; çarpmayı ise ``$\alpha$'nın kopyalarından oluşan bir
$\beta$-dizisi'' diye açıkladık. Fakat
$\text{finfun}(\alpha,\beta)$ hakkında böylesine özlü bir açıklamamız
yoktur. Bu nedenle sıra sayılarında üs almayı bunun yerine doğrudan
sonlu ötesi özyinelemeyle şöyle tanımlayacağız:

\begin{defn}\ollabel{ordexporecursion}
\begin{align*}
	\ordexpo{\alpha}{0} &= 1\\
	\ordexpo{\alpha}{\beta\ordplus 1} &=\ordexpo{\alpha}{\beta} \ordtimes \alpha\\
	\ordexpo{\alpha}{\beta} &= \bigcup_{\delta < \beta}\ordexpo{\alpha}{\delta}& & \text{$\beta$ bir limit sıra sayısı olduğunda}
\end{align*}
\end{defn}

Küme teorisyenleri \emph{sıfatıyla} çalışıyor olsaydık sıra sayılarında üs
almanın bazı özelliklerini araştırmak isteyebilirdik. Burada, üs almanın
değişmeli olmadığına ilişkin şaşırtıcı olmayan olgu dışında ekleyecek fazla
bir şeyimiz yoktur. Nitekim
$\ordexpo{2}{\omega}=\bigcup_{\delta<\omega}\ordexpo{2}{\delta}=\omega$,
oysa $\ordexpo{\omega}{2}=\omega\ordtimes\omega$'dır. Zaten doğal sayılarda
da değişmeli olmadığı için üs almanın değişmeli olmasını
\emph{beklememeliyiz}: $\ordexpo{2}{3}=8<9=\ordexpo{3}{2}$.

\begin{prob}
Sonlu Ötesi Tümevarım'ı kullanarak
$\ordexpo{\alpha}{\beta}=\ordtype{\text{finfun}(\alpha,\beta),\sqsubset}$
diye tanımladığımızda
\olref[sth][ord-arithmetic][expo]{ordexporecursion} içindeki özyineleme
denklemlerini elde ettiğimizi ispatlayın.
\end{prob}

